\begin{center}
\thispagestyle{empty}
\vspace{2cm}
\LARGE{\textbf{RESUMEN}}\\[1.0cm]
\end{center}
\thispagestyle{empty}

Este informe presenta el avance del proyecto “Lector PDF Minimalista”, una aplicación de escritorio orientada a la lectura y gestión de documentos PDF. El proyecto surge como respuesta a la necesidad de contar con una herramienta sencilla, organizada y centrada en el contenido, evitando la sobrecarga de botones, menús y funciones que caracteriza a algunos lectores disponibles actualmente. La aplicación está siendo desarrollada principalmente para Windows utilizando Electron, JavaScript, HTML, CSS y la biblioteca PDF.js.

Durante el segundo sprint se construyó TODO.

Además, el informe documenta los requerimientos iniciales del sistema, el producto mínimo viable correspondiente al Sprint 2, la organización del equipo, la arquitectura propuesta, los prototipos de interfaz y el proceso de control de calidad. Esta entrega representa una continuación sobre la base funcional desarrollada en etapas anteriores, que inluye las anotaciones, los comentarios, la firma gráfica y la integración con el sistema de firma digital GAUDI, así como el un instalador para configurar el programa en el sistema operativo. \\

\textbf{Palabras clave: }lector PDF, Electron, PDF.js, interfaz minimalista, firma digital, comentarios, highlights, GAUDI.